\section{Divisor volumes and the plan of proof}
\label{paper:arithmetic-setup}

The two counting problems have a common geometric representation.
For $3\leq y_1\leq\cdots\leq y_k$, let $\mathcal A_k(\mathbf y)$
consist of squarefree tuples $\mathbf a=(a_1,\ldots,a_k)$ with prime
factors of $a_1$ at most $y_1$, and those of $a_i$, $i\geq2$, in
$(y_{i-1},y_i]$. Thus $2$ and $3$ belong to the first prime range.
The components are pairwise coprime, and an empty selection in any
component is allowed. Define
\begin{align}
\mathcal L^{(k+1)}(\mathbf a)
 &=\bigcup_{\substack{d_1,\ldots,d_k\in\mathbb N\\
       d_1\cdots d_j\mid a_1\cdots a_j\ (1\leq j\leq k)}}
       \prod_{j=1}^k[\log(d_j/2),\log d_j),\label{paper:divisor-boxes}\\
L^{(k+1)}(\mathbf a)&=\operatorname{vol}_k\mathcal L^{(k+1)}(\mathbf a),\notag\\
K_k(\ell)&=e^{-\sum_{i=1}^k h_iL_i}
       \sum_{\mathbf a\in\mathcal A_k(\mathbf y)}
       \frac{L^{(k+1)}(\mathbf a)}{a_1\cdots a_k},
       \qquad h_i=k-i+2.\label{paper:K}
\end{align}
Here $L_i=\ell_i-\log3\geq0$ and
$\log y_i=e^{L_i}\log y_{i-1}$, starting from $y_0=3$.
The volume $L^{(k+1)}$ and the block lengths $L_i$ are different objects.
The sum in \eqref{paper:K} includes all such squarefree tuples.
The empty tuple alone contributes $(\log2)^k$ before normalization.

A prime from the $i$th range can be assigned to one of the divisor
coordinates $i,\ldots,k$, or left unused. These are its $h_i$ eligible
labels. The union in \eqref{paper:divisor-boxes} records all the
resulting decompositions, counting overlapping boxes only once.
This overlap, rather than the number of decompositions, is the source
of the difficulty. The harmonic weight makes it possible to study the
prime locations and their labels separately while retaining the original
divisor union.

Koukoulopoulos's arithmetic comparison
\cite[Theorem~1.7]{Kouk} gives
\begin{equation}
\frac{H^{(k+1)}(x,\mathbf y,2\mathbf y)}x
 \asymp_k K_k(\ell(\mathbf y))
\quad\text{if}\quad
x\geq 2^k y_1\cdots y_k y_k.
\label{main:admissible}
\end{equation}
The first prime range in that comparison begins below $y_1$ and includes
the small primes; the choice $y_0=3$ above fixes only the logarithmic
normalization. Bounded changes in these normalizing factors are absorbed
by constants depending on $k$. We prove the table transfer, including
bounded axes, in Section~\ref{paper:assembly}.

\begin{theorem}\label{main:uniform}
For every fixed $k\geq6$, the function defined in
\eqref{profile:final} satisfies
\[
 K_k(\ell)\asymp_k\Phi_k(\ell)
 \qquad(\ell\in[\log3,\infty)^k).
\]
\end{theorem}

There are three parts to the proof. On positive physical lengths
$Q_i\geq1$, a covering argument gives
$K_k(\log3+Q)\ll_k\Psi_k(Q)$. A separate construction gives
$K_k(\log3+L)\gg_k\Psi_k(L^{L_*})$, where
$L_i^{L_*}=\max(L_i,L_*)$ and $L_*=L_*(k)$ is fixed once.
Finally, stability of the original harmonic sum compares the arithmetic
quantities at $L$ and $L^\lambda$. It allows the unknown threshold
$L_*$ to be removed from the definition of the answer by taking the
explicit infimum in \eqref{profile:final}.

The same ordered saddle appears in both estimates. Its tilt extracts
the exponential cost $e^{-J}$. Within each maximal equality block of
the saddle, a shared prime clock acts as an environment for several
private label constraints. This is why a fractional moment of a
conditional survival probability appears: averaging the environment
before taking the owner power would describe a different event.
Exit-count windows contribute their own powers. The final ordinary
suffix has one count variable that must simultaneously supply the root
and permit path survival. The proof keeps these requirements together.

The deterministic definitions come first, followed by explicit examples
and the saddle geometry. We then establish the scalar persistence and
conditioned bridge estimates, construct the common accepting measure
for the lower bound, and prove the upper bound directly for the full
divisor union. Arithmetic stability closes the proof on the original
parameter domain. We then return to the same lower source to identify
the contributing clocks and interpret the profile's entropy costs in
Section~\ref{fp:sec:arithmetic}.

\subsection{Size truncation of the harmonic volume sum}
\label{artrunc:section}

The arithmetic quantity in the main theorem sums all squarefree tuples
in their prescribed prime blocks. The actual rectangle arguments may
nevertheless start with a size-truncated tuple. The following positive
weighted truncation proves this implication independently of any sharp
upper bound, Poisson approximation, saddle point, or proposed profile.

Fix $H=k+1$ and $3\leq y_1\leq\cdots\leq y_k$. Write
$V(\mathbf a)=L^{(k+1)}(\mathbf a)$ and
\[
 S(\mathbf y)=\sum_{\mathbf a\in\mathcal A_k(\mathbf y)}
             \frac{V(\mathbf a)}{a_1\cdots a_k}.
\]
All prime sets here are finite; in particular all following rearrangements
are identities between finite nonnegative sums. Choose one absolute
constant $C_{\rm W}$ for the elementary weighted prime estimate
\begin{equation}
 \sum_{p\leq y}\frac{\log p}{p}\leq C_{\rm W}\log y
 \qquad(y\geq3).
 \label{artrunc:weighted-primes}
\end{equation}
For example the usual Chebyshev bound for $\sum_{p\leq y}\log p$ and
partial summation give this estimate; see \cite{RosserSchoenfeld1962}. Put
\[
 \tau_k=2kH C_{\rm W}+1,
 \qquad
 S_{\rm tr}(\mathbf y)=
 \sum_{\substack{\mathbf a\in\mathcal A_k(\mathbf y)\\
                         a_i\leq y_i^{\tau_k}\ (1\leq i\leq k)}}
                \frac{V(\mathbf a)}{a_1\cdots a_k}.
\]

\begin{lemma}[Uniform positive weighted truncation]
\label{artrunc:lemma}
Uniformly in the ordered cutoffs,
\begin{equation}
 S_{\rm tr}(\mathbf y)\leq S(\mathbf y)
                   \leq2 S_{\rm tr}(\mathbf y).
 \label{artrunc:comparison}
\end{equation}
The same comparison holds after multiplication by the exact normalizing
factor defining $K_k$. Every truncated tuple has, for each $i$,
\begin{equation}
 \log(a_1\cdots a_i)
       \leq\tau_k\sum_{g\leq i}\log y_g
       \leq k\tau_k\log y_i.
 \label{artrunc:prefix-size}
\end{equation}
\end{lemma}

\begin{proof}
Let $p$ be a prime in original block $i$, and remove it from that
coordinate of a squarefree tuple. Inserting $p$ back permits at most
$k-i+2\leq H$ eligible routing labels, including the unused label.
For each label the new divisor-box union is a translate of the old
one. Therefore, if $p$ does not already divide the tuple,
\begin{equation}
 V(b_1,\ldots,pb_i,\ldots,b_k)\leq H V(\mathbf b).
 \label{artrunc:slot-growth}
\end{equation}
Using the squarefree identity $\log a_i=\sum_{p\mid a_i}\log p$ gives
the exact decomposition
\begin{align}
 \sum_{\mathbf a\in\mathcal A_k(\mathbf y)}
       \frac{V(\mathbf a)\log a_i}{a_1\cdots a_k}
 &=\sum_{p\text{ in block }i}\frac{\log p}{p}
       \sum_{\substack{\mathbf b\in\mathcal A_k(\mathbf y)\\p\nmid b_i}}
       \frac{V(b_1,\ldots,pb_i,\ldots,b_k)}{b_1\cdots b_k}
       \notag\\
 &\leq H S(\mathbf y)
                \sum_{p\text{ in block }i}\frac{\log p}{p}
 \leq H C_{\rm W}\log y_i\,S(\mathbf y).
 \label{artrunc:first-moment}
\end{align}
The prime blocks are disjoint, so removal and reinsertion in this
identity preserve every original prime-range restriction. In the
inequality we only enlarged the inner nonnegative sum.

For the bad set $a_i>y_i^{\tau_k}$, multiply its indicator by
$\log a_i/(\tau_k\log y_i)\geq1$ and use
\eqref{artrunc:first-moment}. A union bound over $i$ yields
\[
 S(\mathbf y)-S_{\rm tr}(\mathbf y)
 \leq\frac{kH C_{\rm W}}{\tau_k}S(\mathbf y)
 <\frac12S(\mathbf y).
\]
This proves \eqref{artrunc:comparison}. The common normalizing factor
does not change its ratio. Finally \eqref{artrunc:prefix-size} follows
from the individual size bounds and the ordering of the cutoffs.
\end{proof}

In particular, every divisor coordinate whose available primes lie in
the first $i$ original groups has logarithmic width at most
$C_k\log y_i$ on the truncated tuples. Earlier contributions to a
remaining coordinate obey the same bound at the last earlier cutoff.
These are precisely the full-size baseline rectangles used by the
actual routing covers. The constant $C_k=k\tau_k$, and its fixed
powers, are absorbed into their dimension constants and fixed reserves,
with the physical lengths, KKT products and owners fixed. Apply those covers
first to this truncated sum, and only then perform the positive repeated-
prime and cell enlargements. Lemma~\ref{artrunc:lemma} supplies the upper
bound for the original full sum. The size inequalities enter at the
truncated squarefree stage, before the repeated-prime enlargement.
